\section{Conclusion}
\label{sec:conclusion}

We have traced the behavior of a spin-$1/2$ Heisenberg pentamer chain from
isolated local multiplets to collective response.
At the tested positive couplings, the pentamer retains an even-parity spin-$1/2$
doublet, positive inter-cell exchange generates an antiferromagnetic effective
chain, and the leading projection error is quadratic in the inter-cell coupling.
Entanglement, finite-size gaps, correlations, and static structure factors
identify a $c=1$ effective chain at those points.
Reflection parity makes the lowest excited doublet dark at $q=\pi$ and shifts most
local inelastic weight into spin-$3/2$ sectors.
Complete connected real-time spectra at $(\alpha,\beta)=(\pm0.5,0.5)$ show
different distributions of low- and high-energy weight in open $L=40$ chains.
Their common finite-time width is $0.09277J_1$, and a joint refinement supports
stability at five tested momenta.
The earlier correction-vector comparison is restricted to a different parameter
point and real-time estimator; the collective bands have not yet been assigned
quantitatively to individual local multiplets.

On the negative-$\alpha$ side, the isolated cell switches to a spin-$3/2$ quartet
whose shape factor favors the $0.60\pi$ replica among folded $Q=\pi$ momenta
at every sampled point, and the
accepted calculations reveal a reconstruction from $\qpeak\simeq0.60\pi$ to
$q=\pi$.
All 29 final refinement calculations converge, and a
terminal $k=8\rightarrow12$ control shows that the most sensitive $L=80$ endpoint
is stable with respect to a larger Krylov subspace.
This endpoint control does not establish solver stability of the entire size
sequence or scan.
None of the eight candidate intervals satisfies the momentum and dominant-ratio
criteria simultaneously, so the reconstruction remains a finite-size result
compatible with a crossover rather than an established thermodynamic boundary.

For an odd-unit-cell cluster chain, the low-energy physics is controlled not by an
inherited phase label but by which local multiplet survives, how it couples
between cells, and which of its excitations the symmetry of the probe allows one
to see.
